\section{KnowNo Baseline Details}~\label{appx:exp}

This appendix describes the KnowNo-style action-query baseline used in the policy comparison. The baseline uses the same domain scenarios, task dynamics, execution limits, and scene randomization protocol as the other policy baselines. The main experiments report the GPT-4o variant with the structured prompt version shown below, conformal threshold $\hat{q}=0.92$, and scoring temperature $3.0$. This setup follows the action-query framing of KnowNo while keeping the environment, transition model, and maximum step budget matched to the state-query baselines.

\subsection{KnowNo Baseline Implementation}~\label{appx:knowno_impl}

\paragraph{Prompt structure.}
At each decision step, KnowNo first generates multiple-choice action options and then scores those options. The GPT-4o and GPT-3.5-turbo variants use the same prompt templates, the same runtime state fields, the same conformal threshold $\hat{q}=0.92$, and the same scoring temperature $3.0$; only the underlying language model is changed. Listings~\ref{lst:knowno_prompt_tomato} and~\ref{lst:knowno_prompt_waste} show the prompt structure used in the experiments. The actual generation prompt includes domain-specific few-shot examples before the runtime state block.

\begin{tcblisting}{
breakable,
listing only,
title=KnowNo tomato-harvesting prompt template,
label={lst:knowno_prompt_tomato},
title filled=false,
colback=black!5!white,
colframe=black!75!black,
left=2pt,
right=2pt,
top=2pt,
bottom=2pt,
listing options={
basicstyle=\scriptsize\ttfamily,
breaklines=true,
columns=fullflexible
}}
Generation prompt
-----------------
We: You are a tomato harvesting robot. The robot moves between a dock station and tomato stems, observes tomatoes, picks ripe tomatoes, scans held tomatoes for quality, loads good tomatoes, and discards rotten tomatoes.

Action roles:
- navigate to <location>: move the robot to dock_station, stem_01, or stem_02.
- detect <location>: observe tomatoes at the current robot stem.
- pick <tomato>: pick one detected ripe tomato at the current robot stem with an empty hand.
- scan <tomato>: inspect the currently held tomato.
- place <tomato>: load the held ripe tomato.
- discard <tomato>: discard the held rotten or bad tomato.

Output rules:
- Each option must be exactly one allowed action from the action roles.
- Use only these action formats: navigate to <location>, detect <location>, pick <tomato>, scan, scan <tomato>, place <tomato>, discard <tomato>.
- Do not generate descriptive phrases, retries, requests for assistance, system checks, or any action outside these formats.

<domain-specific few-shot examples for valid option generation>

We: The task is: <instruction>
We: The environment contains four tomatoes: tomato1, tomato2, tomato3, and tomato4. The locations are dock_station, stem_01, and stem_02.
We: The robot is at <robot_location>. <held_tomato_sentence>
We: The tomatoes that still need attention are <active_tomatoes>.
We: <tomato_state_sentences>
We: The loaded tomatoes are <loaded_tomatoes>. The discarded tomatoes are <discarded_tomatoes>.
We: <completed_action_history>
We: What should the robot do next? Answer with four options labeled A), B), C), and D).
You:

Scoring prompt
--------------
You are a tomato harvesting robot. The robot moves between a dock station and tomato stems, observes tomatoes, picks ripe tomatoes, scans held tomatoes for quality, loads good tomatoes, and discards rotten tomatoes.

Action roles:
- navigate to <location>: move the robot to dock_station, stem_01, or stem_02.
- detect <location>: observe tomatoes at the current robot stem.
- pick <tomato>: pick one detected ripe tomato at the current robot stem with an empty hand.
- scan <tomato>: inspect the currently held tomato.
- place <tomato>: load the held ripe tomato.
- discard <tomato>: discard the held rotten or bad tomato.

We: The task is: <instruction>
We: The environment contains four tomatoes: tomato1, tomato2, tomato3, and tomato4. The locations are dock_station, stem_01, and stem_02.
We: The robot is at <robot_location>. <held_tomato_sentence>
We: The tomatoes that still need attention are <active_tomatoes>.
We: <tomato_state_sentences>
We: The loaded tomatoes are <loaded_tomatoes>. The discarded tomatoes are <discarded_tomatoes>.
We: <completed_action_history>
We: What should the robot do next?
You:
<generated_options>
We: Which option is correct? Answer with a single capital letter.
You:
\end{tcblisting}

\begin{tcblisting}{
breakable,
listing only,
title=KnowNo waste-sorting prompt template,
label={lst:knowno_prompt_waste},
title filled=false,
colback=black!5!white,
colframe=black!75!black,
left=2pt,
right=2pt,
top=2pt,
bottom=2pt,
listing options={
basicstyle=\scriptsize\ttfamily,
breaklines=true,
columns=fullflexible
}}
Generation prompt
-----------------
We: You are a robot operating in a waste sorting station. You are in front of a counter. There are four bins: a general bin, a plastic bin, a paper bin, and a can bin.

Action roles:
- detect: observe waste attributes for objects on the counter.
- pick <object>: pick one observed object from the counter with an empty hand.
- place <object> into <bin>: place the held object into the bin matching the observed attribute.

Output rules:
- Each option must be exactly one allowed action from the action roles.
- Use only these action formats: detect, pick <object>, place <object> into <bin>.
- Do not generate descriptive phrases, retries, requests for assistance, system checks, or any action outside these formats.

<domain-specific few-shot examples for valid option generation>

We: The task is: <instruction>
We: The waste objects still on the counter are <remaining_objects>.
We: The available bins are <available_bins>.
We: The occluded waste objects are <occluded_objects>.
We: <observed_waste_sentence> <held_object_sentence>
We: <completed_action_history>
We: What should the robot do next? Answer with four options labeled A), B), C), and D).
You:

Scoring prompt
--------------
You are a robot operating in a waste sorting station. You are in front of a counter. There are four bins: a general bin, a plastic bin, a paper bin, and a can bin.

We: The task is: <instruction>
We: The waste objects still on the counter are <remaining_objects>.
We: The available bins are <available_bins>.
We: The occluded waste objects are <occluded_objects>.
We: <observed_waste_sentence> <held_object_sentence>
We: <completed_action_history>
We: What should the robot do next?
You:
<generated_options>
We: Which option is correct? Answer with a single capital letter.
You:
\end{tcblisting}

\paragraph{Prediction set and $\hat{q}$.}
For each option letter, the implementation extracts the model's top-token log probability and applies temperature scaling with temperature $3.0$ to obtain normalized option scores. The prediction set is then
\[
\Gamma_{\hat{q}}(x)=\{y \in \mathcal{Y}(x) \mid score(y \mid x) \ge 1-\hat{q}\},
\]
where $x$ denotes the current prompt context and $\mathcal{Y}(x)$ is the set of generated option letters. If $\Gamma_{\hat{q}}(x)$ contains a single executable option, KnowNo executes that option. If the set contains multiple options or includes a fallback/ambiguous option, the baseline requests user help at the action level.

Statistically, $\hat{q}$ is the conformal prediction-set threshold. It is computed once from a calibration set by measuring the nonconformity of the correct option,
\[
r_i = 1 - \max_{y \in Y_i^{correct}} score(y \mid x_i),
\]
and taking the finite-sample conformal quantile
\[
\hat{q} = \mathrm{Quantile}_{\lceil (n+1)(1-\alpha) \rceil/n}\left(\{r_i\}_{i=1}^{n}\right).
\]
Under the usual exchangeability assumption for calibration and test examples, this threshold gives marginal coverage for the correct option at the target level. In practical terms, a larger $\hat{q}$ lowers the score threshold $1-\hat{q}$ and produces larger prediction sets, which increases the chance of asking for help; a smaller $\hat{q}$ produces smaller sets but may omit the correct action more often. In the reported experiments, $\hat{q}=0.92$ is fixed during execution rather than updated online.

\paragraph{Model variant comparison.}
The main policy comparison reports the GPT-4o KnowNo baseline because it was the stronger and more stable action-query variant. We also evaluated the same KnowNo pipeline with GPT-3.5-turbo using the same structured prompt version, $\hat{q}=0.92$, and scoring temperature $3.0$. Table~\ref{tab:app_knowno_model_compare} summarizes the model comparison. GPT-4o achieved higher success with substantially fewer action-level queries and smaller prediction sets. GPT-3.5-turbo produced larger prediction sets and more frequent help requests, and it was especially unstable in tomato harvesting, where many runs terminated after non-executable or invalid action choices. This comparison indicates that the action-query baseline is sensitive to the underlying language model; therefore, the main text uses GPT-4o as the representative KnowNo condition.

\begin{table*}[t]
\centering
\caption{KnowNo model variant comparison. Values are means over logged runs. GPT-4o uses $200$ runs per domain. GPT-3.5-turbo uses $200$ waste-sorting runs and $69$ tomato-harvesting runs available in the experiment logs.}
\label{tab:app_knowno_model_compare}
\small
\setlength{\tabcolsep}{4pt}
\begin{tabular*}{\textwidth}{@{\extracolsep{\fill}}llcccccc}
\toprule
\textbf{Model} & \textbf{Domain} & \textbf{$n$} & \textbf{Success (\%)} & \textbf{Steps} & \textbf{Queries} & \textbf{Pred. set size} & \textbf{Time (s)} \\
\midrule
GPT-4o & Waste & 200 & 38.5 & 7.23 & 2.06 & 1.95 & 23.68 \\
GPT-4o & Tomato & 200 & 42.0 & 10.73 & 2.12 & 2.17 & 43.17 \\
GPT-4o & All & 400 & 40.2 & 8.98 & 2.10 & 2.06 & 33.43 \\
\midrule
GPT-3.5-turbo & Waste & 200 & 34.5 & 6.42 & 7.00 & 3.44 & 47.55 \\
GPT-3.5-turbo & Tomato & 69 & 2.9 & 5.14 & 6.07 & 3.51 & 45.24 \\
GPT-3.5-turbo & All & 269 & 26.4 & 6.09 & 6.76 & 3.46 & 46.96 \\
\bottomrule
\end{tabular*}
\end{table*}
